\chapter[Highest Weights and Finite-Dimensional Modules]{Highest Weights\\and Finite-Dimensional Modules}
\label{ch:highest-weights}

In this chapter \(\mathfrak g\) is a complex semisimple Lie algebra.
Fix a Cartan subalgebra \(\mathfrak h\), positive roots \(R^+\),
simple roots \(\alpha_1,\ldots,\alpha_r\), and normalized triples
\((e_i,f_i,h_i)\). Put
\[
 Q_+=\sum_i\mathbb Z_{\geq0}\alpha_i,\qquad
 P=\{\lambda\in\mathfrak h^*:\lambda(h_i)\in\mathbb Z\},\qquad
 P_+=\{\lambda\in P:\lambda(h_i)\geq0\}.
\]
The order \(\mu\leq\lambda\) means \(\lambda-\mu\in Q_+\).
The weight space in a module \(V\) is
\(V_\mu=\{v:hv=\mu(h)v\text{ for all }h\in\mathfrak h\}\).
The root-space results give
\(\mathfrak g_\alpha V_\mu\subseteq V_{\mu+\alpha}\).
Our arguments do not use the integration of a representation to a
Lie group. The sources are \cite[Sections 20--21]{humphreys} and
\cite[Section 25]{etingof}.

\section{Induction and the universal highest-weight module}

\begin{proposition}[Triangular multiplication and Verma modules]
\label{res:G01}
Multiplication is a vector-space isomorphism
\[
 U(\mathfrak n^-)\otimes U(\mathfrak h)\otimes U(\mathfrak n^+)
 \longrightarrow U(\mathfrak g).
\]
For \(\lambda\in\mathfrak h^*\), let \(\mathbb C_\lambda\) be the
one-dimensional left module over
\(\mathfrak b=\mathfrak h\oplus\mathfrak n^+\), with \(\mathfrak n^+\)
acting by zero and \(h\) acting by \(\lambda(h)\). Then
\[
 M(\lambda)=U(\mathfrak g)\otimes_{U(\mathfrak b)}\mathbb C_\lambda
 \simeq U(\mathfrak g)/
 \left(\sum_{h\in\mathfrak h}U(\mathfrak g)(h-\lambda(h))
       +\sum_iU(\mathfrak g)e_i\right).
\]
The map \(u\mapsto uv_\lambda\) identifies \(U(\mathfrak n^-)\)
linearly with \(M(\lambda)\), where \(v_\lambda=1\otimes1\).
\end{proposition}
\begin{proof}
Choose ordered bases of the negative root spaces, \(\mathfrak h\),
and the positive root spaces, in that order. PBW
(Theorem~\ref{res:A05}) says exactly that multiplication sends the
tensor product of their ordered monomial bases bijectively to the
ordered monomial basis of \(U(\mathfrak g)\).
In particular \(U(\mathfrak g)\) is free as a right
\(U(\mathfrak b)\)-module on the negative-root monomials.

The prescribed action on \(\mathbb C_\lambda\) is a Lie action:
\([\mathfrak b,\mathfrak b]\subseteq\mathfrak n^+\), so its image
vanishes. The induced associative algebra character has kernel the
left ideal of \(U(\mathfrak b)\) generated by
\(h-\lambda(h)\) and \(e_i\). To check this description, commute all
positive generators to the right in PBW order. Every positive-degree
monomial in the positive generators lies in their generated left ideal;
the remaining Cartan polynomial reduces to its value at \(\lambda\).
The simple \(e_i\) generate \(\mathfrak n^+\), so they suffice.
Tensoring the right-free decomposition with that one-dimensional
quotient proves both identifications and the final assertion.
\end{proof}

\begin{proposition}[Weights and the universal property]
\label{res:G02}
The module \(M(\lambda)\) is a direct sum of finite-dimensional weight
spaces, has weights precisely \(\lambda-Q_+\), and has
\(M(\lambda)_\lambda=\mathbb C v_\lambda\). A vector \(v\) in any
\(\mathfrak g\)-module satisfying
\[
 hv=\lambda(h)v,\qquad e_iv=0\quad(1\leq i\leq r)
\]
determines a unique homomorphism \(M(\lambda)\to V\) taking
\(v_\lambda\) to \(v\). Every module generated by such a nonzero
vector is a quotient of \(M(\lambda)\).
\end{proposition}
\begin{proof}
A negative-root PBW monomial
\(\prod_{\alpha>0}f_\alpha^{a_\alpha}v_\lambda\) has weight
\(\lambda-\sum_{\alpha>0}a_\alpha\alpha\).
If the latter sum has height \(n\), each \(a_\alpha\leq n\);
there are only finitely many choices. This proves finite weight
spaces. Every element of \(Q_+\) can be obtained using the simple
negative root vectors, and their PBW monomials are nonzero, which
proves the assertion about the weight set. The empty monomial is
the sole one of weight \(\lambda\).
The map \(u\mapsto uv\) kills the displayed ideal in
Proposition~\ref{res:G01}, hence factors through \(M(\lambda)\).
Its value on the generator forces uniqueness and its image is
the submodule generated by \(v\).
\end{proof}

\begin{lemma}[Submodules of a weight module]
\label{supp:weight-submodules}
If \(V=\bigoplus_\mu V_\mu\) and \(W\subseteq V\) is
\(\mathfrak h\)-stable, then
\(W=\bigoplus_\mu(W\cap V_\mu)\), and \(V/W\) has the corresponding
quotient weight spaces.
\end{lemma}
\begin{proof}
A vector has finitely many weight components. Choose \(h\in\mathfrak h\)
which distinguishes those finitely many weights. Lagrange interpolation
in the operator \(h\) extracts each component of the vector. If the
vector lies in \(W\), so does each extracted component. This proves
the direct sum assertion. Applying it to representatives in \(V/W\)
proves the quotient assertion.
\end{proof}

\begin{theorem}[The unique simple highest-weight quotient]
\label{res:G03}
For every \(\lambda\in\mathfrak h^*\), \(M(\lambda)\) has a unique
maximal proper submodule and hence a unique simple quotient
\(L(\lambda)\). Every nonzero highest-weight module of weight
\(\lambda\) surjects onto \(L(\lambda)\).
Every simple finite-dimensional \(\mathfrak g\)-module is one of
these quotients; its highest weight determines its isomorphism class.
\end{theorem}
\begin{proof}
A proper submodule of \(M(\lambda)\) has zero intersection with
its top weight line: containing \(v_\lambda\) would make it the
whole module. Lemma~\ref{supp:weight-submodules} therefore implies
that the sum \(J(\lambda)\) of all proper submodules still has zero
top weight space. It is proper, contains every proper submodule,
and is the unique maximal one. If a nonzero highest-weight quotient
has kernel \(K\), then \(K\subseteq J(\lambda)\), giving the asserted
surjection. For a simple quotient this surjection is an isomorphism.

For a finite-dimensional module, Theorem~\ref{res:C07} supplies
the Cartan weight decomposition. Choose a weight maximal for the
order defined by \(R^+\); it exists because there are finitely many
weights. Every positive-root operator kills that weight space.
A nonzero vector in it generates the module when the module is
simple. The universal property applies.
Finally if two highest-weight generators have weights \(\lambda,\mu\),
all weights of the generated module lie both in \(\lambda-Q_+\)
and in \(\mu-Q_+\). Each of \(\lambda,\mu\) is a weight, so
\(\lambda\leq\mu\leq\lambda\), giving \(\lambda=\mu\).
\end{proof}

\section{The finite-dimensionality criterion}

\begin{proposition}[Necessity of dominance and integrality]
\label{res:G04}
If \(L(\lambda)\) is finite dimensional, then \(\lambda\in P_+\).
More generally, every finite-dimensional highest-weight module of
weight \(\lambda\) satisfies this conclusion and
\(f_i^{\lambda(h_i)+1}v_\lambda=0\).
\end{proposition}
\begin{proof}
Restrict to the \(i\)-th \(\mathfrak{sl}_2\).
Its highest vector \(v_\lambda\) has \(h_i\)-weight \(\lambda(h_i)\).
The rank-one classification, Theorem~\ref{res:C01}, forces this
weight to be a nonnegative integer and gives the asserted last
nonzero step of the \(f_i\)-string. This holds for every \(i\).
\end{proof}

\begin{lemma}[Local finiteness from the generator]
\label{supp:local-integrability}
For \(\lambda\in P_+\), the quotient
\[
 Q(\lambda)=M(\lambda)/
     \sum_i U(\mathfrak g)f_i^{\lambda(h_i)+1}v_\lambda
\]
is nonzero. Both \(e_i\) and \(f_i\) act locally nilpotently on it.
Every vector lies in a finite-dimensional module for the
\(i\)-th \(\mathfrak{sl}_2\).
\end{lemma}
\begin{proof}
Write \(m_i=\lambda(h_i)\).
The rank-one commutator gives
\(e_i f_i^{m_i+1}v_\lambda=0\).
For \(j\ne i\), \([e_j,f_i]=0\), so
\(e_j f_i^{m_i+1}v_\lambda=0\) as well.
These vectors are thus highest vectors of weights strictly below
\(\lambda\). Each of their generated submodules has only weights
below its own highest weight. Their sum has no \(\lambda\) weight,
so \(v_\lambda\) survives in \(Q(\lambda)\).

The adjoint operators of \(e_i,f_i\) are nilpotent on
\(\mathfrak g\), by the root strings. Their derivation extensions
are locally nilpotent on \(U(\mathfrak g)\): on a product of finitely
many generators the repeated derivative eventually differentiates
one factor more often than its nilpotence bound. The identity
\[
 x^n(uv)=\sum_{j=0}^n\binom nj(\operatorname{ad}x)^j(u)x^{n-j}v
\]
therefore proves local nilpotence on \(Q(\lambda)\), starting from
\(e_iv_\lambda=0\) and \(f_i^{m_i+1}v_\lambda=0\).

For a weight vector \(v\), PBW for \(\mathfrak{sl}_2\) spans its
generated module by \(f_i^a h_i^b e_i^c v\).
Only finitely many \(c\) occur; each \(e_i^cv\) is a weight vector,
so \(h_i^b\) contributes only scalar factors; and for each of these
finitely many vectors only finitely many \(a\) occur. Thus the span
is finite. A general vector is a finite sum of weight vectors,
which proves the last assertion.
\end{proof}

\begin{theorem}[Existence of finite-dimensional highest-weight modules]
\label{res:G05}
For every \(\lambda\in P_+\), \(Q(\lambda)\) is finite dimensional.
Its weight multiplicities are invariant under the Weyl group.
Consequently \(L(\lambda)\) is finite dimensional.
\end{theorem}
\begin{proof}
Consider a weight module locally finite for the \(i\)-th
\(\mathfrak{sl}_2\). If \(\mu(h_i)=m\geq0\), then
\(f_i^m:V_\mu\to V_{s_i\mu}\) is injective: any nonzero vector
generates a finite-dimensional rank-one module, and the explicit
rank-one decomposition shows that this operator is injective on
its \(m\) weight space. If \(m\leq0\), use \(e_i^{-m}\) instead.
Applying the same argument to \(s_i\mu\) gives equality of
dimensions. Weights occurring here are integral because restriction
to every rank-one algebra is a sum of its finite-dimensional modules.
Since the \(s_i\) generate \(W\), all multiplicities are
\(W\)-invariant.

We prove finiteness explicitly. If a weight \(\mu\) is dominant,
write \(\mu=\sum_i a_i\omega_i\) with \(a_i\in\mathbb Z_{\geq0}\).
Let \(\rho^\vee\) be half the sum of the positive coroots. Each
\(\omega_i(\rho^\vee)>0\): the coefficient of \(\alpha_i^\vee\)
in this positive sum is positive. Since
\(\alpha_i(\rho^\vee)=1\) and \(\lambda-\mu\in Q_+\), we have
\[
 0\leq\sum_i a_i\omega_i(\rho^\vee)
      =\mu(\rho^\vee)\leq\lambda(\rho^\vee).
\]
This bounds every nonnegative integer \(a_i\); thus there are only
finitely many dominant weights below \(\lambda\).
Every integral weight has a Weyl conjugate in the dominant chamber,
and \(W\) is finite. The weight set of \(Q(\lambda)\) is therefore
finite. Its weight spaces are finite as quotients of those of
\(M(\lambda)\). This proves finite dimensionality, without assuming
it in the local-finiteness argument.
The map \(M(\lambda)\to L(\lambda)\) kills the indicated singular
vectors (or use Proposition~\ref{res:G03}); hence \(L(\lambda)\)
is a quotient of \(Q(\lambda)\) and is finite dimensional.
\end{proof}

\begin{theorem}[Classification and defining relations]
\label{res:G06}
The map \(\lambda\mapsto L(\lambda)\) is a bijection from \(P_+\)
to the simple finite-dimensional \(\mathfrak g\)-modules up to
isomorphism. In addition, \(Q(\lambda)=L(\lambda)\). Thus
\(L(\lambda)\) is generated by one vector \(v\) with precisely the
defining relations
\[
 e_iv=0,\quad hv=\lambda(h)v,\quad
 f_i^{\lambda(h_i)+1}v=0 .
\]
Every finite-dimensional module is a direct sum of these modules.
\end{theorem}
\begin{proof}
Necessity, existence and uniqueness follow from
Propositions~\ref{res:G03}--\ref{res:G04} and
Theorem~\ref{res:G05}. To verify the stronger presentation assertion,
apply Weyl complete reducibility (Theorem~\ref{res:C05}) to
\(Q(\lambda)\). Project its cyclic highest vector onto each simple
summand. A nonzero projection is again a highest vector of weight
\(\lambda\), so that summand is \(L(\lambda)\); a zero projection
would contradict cyclic generation of the whole summand.
The \(\lambda\) weight space has dimension one. Exactly one summand
can therefore occur. The presentation follows.
The last assertion is Weyl complete reducibility with the just-proved
classification of its simple summands.
\end{proof}

\section{Strings, duals and examples}

\begin{proposition}[Weight strings and the convex bound]
\label{res:G07}
For a finite-dimensional module, multiplicities are \(W\)-invariant.
For every root \(\alpha\), the subspace
\(\bigoplus_{j\in\mathbb Z}V_{\mu+j\alpha}\) decomposes into rank-one
strings, each symmetric under \(s_\alpha\).
In particular, if \(m=\langle\mu,\alpha^\vee\rangle\geq0\)
and \(V_\mu\ne0\), all \(\mu-j\alpha\), \(0\leq j\leq m\),
are weights. Every weight of \(L(\lambda)\) belongs to
\[
 (\lambda+Q)\cap\operatorname{conv}(W\lambda).
\]
\end{proposition}
\begin{proof}
The direct sum over the \(\alpha\)-coset is stable under the
root triple. Decompose it into simple rank-one modules using
Theorem~\ref{res:C06}; the Cartan directions killed by \(\alpha\)
act by fixed scalars on this sum and do not affect that
decomposition. In each rank-one module the weights are an unbroken
string, with the explicit raising and lowering coefficients of
Theorem~\ref{res:C01}. Reflection exchanges its two ends, with
equal multiplicities. The claims about strings and multiplicities
follow by summing these decompositions.

All weights are in \(\lambda-Q_+\), so are in \(\lambda+Q\).
For the convex bound let \(H\) lie in the real Cartan space.
Move \(H\) by \(W\) to the closed dominant chamber. For any weight
\(\nu\), \((\lambda-\nu)(H)\geq0\) in that chamber.
The \(W\)-invariance already proved therefore implies, for arbitrary
\(H\),
\(\nu(H)\leq\max_{w\in W}(w\lambda)(H)\).
The separating-hyperplane characterization of the convex hull of
a finite set gives the required inclusion. Equivalently that
characterization follows by minimizing distance to the compact
convex polytope: a point outside it is strictly separated by the
vector from its nearest point.
\end{proof}

\begin{proposition}[Dual highest weights and Cartan products]
\label{res:G08}
The dual of \(L(\lambda)\) is \(L(-w_0\lambda)\).
The tensor product \(L(\lambda)\otimes L(\mu)\) contains one copy
of \(L(\lambda+\mu)\); its highest line is generated by
\(v_\lambda\otimes v_\mu\).
More generally the tensor product of the fundamental modules,
with \(\omega_i\) repeated \(\lambda(h_i)\) times, contains
\(L(\lambda)\) as the submodule generated by their top vectors.
\end{proposition}
\begin{proof}
The dual action is \((x\varphi)(v)=-\varphi(xv)\), so the
dual weights are negatives of the original weights with the same
multiplicities. The lowest weight of \(L(\lambda)\) is \(w_0\lambda\):
apply \(w_0\) to the inequalities \(\nu\leq\lambda\), using that
it changes positive roots to negative roots.
That weight has multiplicity one by Weyl invariance.
Consequently the dual has top weight \(-w_0\lambda\);
it is simple since a proper submodule would annihilate a nonzero
proper submodule quotient of \(L(\lambda)\).

In the tensor product all weights are below \(\lambda+\mu\),
and its top space is exactly the tensor product of the two top
lines. That vector is killed by every \(e_i\).
The generated submodule is a finite-dimensional cyclic highest-weight
module, so Theorem~\ref{res:G06} and its complete-reducibility
argument show that it is \(L(\lambda+\mu)\).
No second copy can occur because the top line is one dimensional.
The assertion about fundamental modules follows by iteration;
for \(\lambda=0\) use the empty tensor product \(\mathbb C\).
\end{proof}

\begin{example}[Type \(A\) and the polynomial comparison]
\label{res:G09}
For \(\mathfrak{sl}_2\), \(L(n\omega_1)=\operatorname{Sym}^n\mathbb C^2\).
For \(\mathfrak{sl}_d\), the defining module has highest weight
\(\omega_1\), \(\bigwedge^k\mathbb C^d\) has highest weight
\(\omega_k\) for \(1\leq k<d\), and
\(\operatorname{Sym}^a\mathbb C^d=L(a\omega_1)\).
For \(\mathfrak{sl}_3\),
\[
 \mathbb C^3\otimes(\mathbb C^3)^*
   =\mathbb C\oplus L(\omega_1+\omega_2),
 \qquad \dim L(\omega_1+\omega_2)=8 .
\]
A polynomial \(\operatorname{GL}_d\)-module indexed by a partition
\(\tau=(\tau_1\geq\cdots\geq\tau_d\geq0)\) restricts to the
\(\mathfrak{sl}_d\) highest weight
\(\sum_{i<d}(\tau_i-\tau_{i+1})\omega_i\).
Two such parameters have the same restricted highest weight exactly
when they differ by an integral power of the determinant.
\end{example}
\begin{proof}
In the defining matrix action \(e_i=E_{i,i+1}\) and
\(h_i=E_{ii}-E_{i+1,i+1}\). On polynomials,
\(e_i=x_i\partial/\partial x_{i+1}\).
A homogeneous polynomial killed by all \(e_i\) is a scalar
multiple of \(x_1^a\): each derivative with respect to
\(x_2,\ldots,x_d\) must vanish in characteristic zero.
Thus a completely reducible symmetric power has a unique highest
line and is simple with the asserted weight.

The weights in the exterior power are the distinct sums of \(k\)
coordinate weights. A basis wedge is killed by all \(e_i\) exactly
when its indices are \(1,\ldots,k\). There cannot be a cancellation
of their images between distinct weights. Therefore it also has
a unique highest line and is simple; evaluating \(h_i\) gives
\(\omega_k\).
The tensor product with the dual is the conjugation representation
on \(\operatorname{End}(\mathbb C^3)\).
Trace splits it into scalar and trace-zero matrices.
The latter is the adjoint module of the simple algebra
\(\mathfrak{sl}_3\); its highest root is
\(\alpha_1+\alpha_2=\omega_1+\omega_2\), and its dimension is eight.

For the final comparison, the polynomial module construction and
its simple highest line are those of
\cite[Chapter 4]{finite-representations}. Differentiation evaluates
the diagonal highest character \(t_1^{\tau_1}\cdots t_d^{\tau_d}\)
on \(h_i\) as \(\tau_i-\tau_{i+1}\).
A Lie-stable subspace is group-stable here: the elementary
unipotent matrices act by exponentials of \(E_{ij}\), and diagonal
matrices act on the same integral weight spaces.
These matrices generate \(\operatorname{GL}_d(\mathbb C)\).
Thus a simple polynomial group module remains simple on
restriction to the Lie algebra after removing its scalar central
action. Finally equality of all successive differences is
equivalent to \(\tau_i-\sigma_i\) being a constant integer for
all \(i\). Multiplication by the corresponding determinant
character has exactly this effect.
\end{proof}
